A la luz de la investigación realizada, elaborar una opinión fundamentada.

A partir de la investigación, puedo darme cuenta de que Kali Linux es una distribución sumamente especializada en el área de ciberseguridad, la cual facilita las tareas de un especialista al proveerle un amplio ecosistema de herramientas preconfiguradas para auditar y mejorar la seguridad de los sistemas informáticos. Considero que tener una disponibilidad tan vasta de herramientas es tan útil como demandante; no solo por el posible uso indebido por parte de actores maliciosos, sino porque puede resultar abrumador para usuarios principiantes debido a la curva de aprendizaje que exige cada utilidad.

Como se observó al analizar las alternativas, es viable utilizar otras distribuciones o sistemas operativos instalando manualmente el software requerido; sin embargo, utilizar Kali Linux durante la formación académica en seguridad informática simplifica significativamente el entorno de trabajo, reduciendo problemas de compatibilidad y dependencias en el despliegue de herramientas especializadas.
